%% =============================================================================
%% 03_SECTION_DRAFT_SKETCHES.tex  --  SAMPLE SKETCHES ONLY. DO NOT INSERT.
%% =============================================================================
%% Purpose: show the intended TONE, STRUCTURE and FLOW (modelled on the
%% high-scoring example report) for each missing section, using THIS thesis's
%% real content. These are opening-paragraph sketches plus bullet notes, NOT
%% finished sections and NOT to be pasted into OVERLEAF_FLAT_BUNDLE/main.tex.
%%
%% Rules honoured here:
%%   - no invented numbers (illustrative numbers are the verified ones from
%%     Chapter 5 / docs/VERIFICATION_SHEET.txt);
%%   - \citep keys used are ALREADY in references.bib (no new references added);
%%   - no universal-ranking claim; caveats kept adjacent; "motivated, not
%%     validated" clinical framing.
%% Wording is original; the example report is a structural model only.
%% =============================================================================


%% ----------------------------------------------------------------------------
%% ABSTRACT  (shape sketch -- one paragraph, ~250-350 words in final form)
%% ----------------------------------------------------------------------------
% Sentence-order template (fill to final length later):
%   [1 problem] Federated learning lets hospitals train a shared model without
%     pooling patient data, but clinical image datasets are severely
%     class-imbalanced and real federations are heterogeneous in both data and
%     compute.
%   [2 approach] This thesis compares FedAvg, FedProx and FedNova on DermaMNIST,
%     a seven-class dermatology benchmark, across controlled statistical- and
%     system-heterogeneity regimes, reporting test macro-F1 at the
%     best-validation checkpoint.
%   [3 headline numbers] Under Li-style straggler asymmetry the FedProx advantage
%     (+0.115 macro-F1) is mediated almost entirely by gamma-inexact
%     partial-update handling (+0.111) rather than the proximal term (+0.004) ...
%   [4 second result] ... and FedNova is robust under deterministic
%     learning-rate asymmetry yet collapses under random-tau stragglers.
%   [5 framing + caveat in one breath] The contribution is a regime map of when
%     each method preserves or loses rare-class signal -- not a universal ranking
%     -- established on a single low-resolution dataset with simulated stragglers
%     and n=3 mechanism probes.
% NOTE: write the Abstract LAST.


%% ----------------------------------------------------------------------------
\section*{1 Introduction (opening sketch)}
% Model opening paragraph:
Training accurate diagnostic models increasingly depends on data that cannot
leave the institution that holds it. Federated learning addresses this by keeping
patient images on-site and sharing only model updates, but it inherits two
difficulties that are especially acute in medicine: the data held by different
hospitals is not identically distributed, and the institutions that hold the
rarest, most clinically important cases are not guaranteed to contribute equally
to the shared model. On a severely imbalanced benchmark such as DermaMNIST, a
model can post high overall accuracy while quietly failing on exactly the rare
lesions that matter, so both the optimiser used to combine client updates and the
metric used to judge them are decisions with clinical consequences.
% Rest of section should cover:
%  - why optimiser choice (not just architecture) is the lever under study;
%  - the gap: comparisons rarely isolate *why* an optimiser wins under a given
%    heterogeneity, and rarely use a rare-class-sensitive metric on a medical task;
%  - a forward pointer to the regime-map contribution.
% Transition into Motivation:
%  "We begin by setting out why each of these pressures -- privacy, imbalance,
%   and heterogeneity -- bears on the choice of federated optimiser."


%% ----------------------------------------------------------------------------
\subsection*{1.1 Motivation (opening sketch)}
Medical image collections are fragmented across hospitals by privacy regulation
and governance, and pooling them centrally is often impossible; federated
learning is the standard response \citep{rieke2020future,sheller2020medicalfl}.
Two properties of real federations make this harder than the centralised case.
First, class prevalence varies sharply between sites, and the dermatology data
studied here is dominated by a single benign class while the diagnostically
urgent lesions are rare. Second, sites differ in compute and availability, so
some clients return only partial work in a given round. Both pressures act on the
same vulnerable quantity -- the contribution of whichever client holds the rare
classes -- which motivates studying the optimiser as the lever that either
protects or discards that contribution.
% Rest: clinical stakes of rare classes (melanoma); why accuracy is the wrong
%   metric here; one sentence each, no results.
% Transition: "These pressures raise four concrete questions."


%% ----------------------------------------------------------------------------
\subsection*{1.3 Research questions (sketch)}
% Short framing sentence then enumerated list (answers deferred):
This thesis asks, for imbalanced medical federated learning, \emph{when} the
choice of optimiser and protocol preserves rare-class performance:
\begin{itemize}
  \item \textbf{RQ1.} Does statistical heterogeneity alone separate FedAvg and
        FedProx?
  \item \textbf{RQ2.} Under straggler asymmetry, is FedProx's advantage carried
        by the proximal term or by gamma-inexact partial-update handling?
  \item \textbf{RQ3.} Is FedNova's robustness regime-dependent -- holding under
        deterministic learning-rate asymmetry but not under random-$\tau$
        stragglers?
  \item \textbf{RQ4.} Is the FedProx advantage a disguised class-imbalance
        correction, or a distinct drift/protocol intervention?
\end{itemize}
% Transition: "Chapter 5 answers each in turn; we state our contributions next."


%% ----------------------------------------------------------------------------
\subsection*{1.4 Contributions (sketch)}
% Lead sentence + enumerated, evidence-tied, hedged:
The thesis delivers a regime map rather than a ranking. Concretely, it
\begin{enumerate}
  \item maps when FedAvg, FedProx and FedNova preserve or lose rare-class signal
        on an imbalanced medical benchmark;
  \item decomposes the FedProx straggler-regime advantage into a proximal-term
        and a gamma-inexact partial-update component, with a matched negative
        control that fixes when the effect can appear;
  \item identifies FedNova's regime-dependence across deterministic and random
        system heterogeneity;
  \item unifies the experiments under a single rare-client signal-flow rule and
        identifies rare-to-mel-nevi collapse as the shared clinical failure mode;
  \item falsifies the reading of FedProx as a mere class-imbalance corrector via
        a loss-side comparator.
\end{enumerate}
% Transition: "Section 1.5 previews how the remaining chapters develop these."


%% ----------------------------------------------------------------------------
\section*{2 Background and Related Work (opening sketch)}
Federated optimisation begins with FedAvg \citep{mcmahan2017fedavg}, which
averages client models in proportion to their data, and the methods studied here
are best understood as responses to its known weaknesses under heterogeneity.
This chapter introduces those weaknesses in the order the thesis later probes
them: statistical heterogeneity in how classes are distributed across clients,
system heterogeneity in how much work each client contributes, and the
optimiser- and loss-level mechanisms proposed to cope with each. Throughout, we
keep the proximal term of FedProx and its gamma-inexact handling of partial
updates conceptually distinct, because separating them is central to the thesis.
% Rest: the 2.1-2.10 sections per the writing plan; each ends by naming what
%   prior work leaves unresolved.
% Transition into the gap paragraph below.


%% ----------------------------------------------------------------------------
\subsection*{2.10 Gap in prior work (model gap paragraph)}
Three questions remain underexplored in this literature. First, FedProx is
typically evaluated as a single method, so it is rarely established whether its
benefit under stragglers comes from the proximal regulariser or from accepting
partial updates -- two mechanisms the original formulation bundles together
\citep{li2020fedprox}. Second, FedNova's normalisation is analysed under the
assumption that per-client work is stable \citep{wang2020fednova}, leaving its
behaviour under randomly varying work largely uncharacterised on imbalanced
medical data. Third, optimiser-level and loss-level responses to imbalance are
seldom compared on the same medical benchmark with a rare-class-sensitive metric.
This thesis addresses the three together by mapping the regimes in which each
choice preserves rare-class signal.
% Transition: "We formalise this setting in Chapter 3."


%% ----------------------------------------------------------------------------
\section*{3 Problem Formulation (opening sketch)}
We consider a federation of $K$ clients, each holding a private training set
drawn from its own class distribution, that cooperate over $R$ rounds to learn a
single global model without exchanging data. In each round the server broadcasts
the current model, every participating client takes some number of local
optimisation steps, and the server combines the returned updates. The three
optimisers studied differ only in how the local objective is formed and how the
returned updates are combined; the heterogeneity we impose differs only in which
clients hold which classes and in how much local work each client completes. The
remainder of this chapter defines these quantities precisely and names the single
concept -- rare-client signal flow -- around which the results are organised.
% Rest: 3.1-3.6 per writing plan; introduce notation at point of use; may \ref
%   the L1/L4 schematic (fig:l1-l4-schematic) in 3.4 (also fixes the unreferenced
%   figure).
% Transition: "Chapter 4 instantiates this formulation as concrete experiments."


%% ----------------------------------------------------------------------------
\section*{6 Discussion (opening sketch)}
The results in Chapter~5 are best read not as a contest between three optimisers
but as a map of the conditions under which each one keeps rare-class signal alive
in the global model. Where the client holding the rare classes continues to
contribute to aggregation, rare-class performance is preserved; where that
contribution is dropped, under-weighted, or drowned in noise, it collapses --
and the same failure surfaces clinically as rare lesions being routed into the
dominant mel-nevi class. This chapter develops that reading, relates it to the
FedAvg, FedProx and FedNova literature, draws cautious implications for
cross-hospital deployment, and states plainly what the evidence does not support.
% Rest: 6.1-6.8 per writing plan; NO new numbers; hedge comparisons.
% Transition between 6.2 and 6.3: "If signal flow is the operative variable, the
%   practical question becomes which heterogeneity a given deployment expects."


%% ----------------------------------------------------------------------------
\subsection*{6.8 Threats to validity (opening sketch)}
Several limitations bound these conclusions, and we state them as consequences of
the design rather than as afterthoughts. Internally, the mechanism probes use
three matched seeds and report directional effects without significance tests;
stragglers are simulated through local-step counts rather than wall-clock
deadlines; and the two runtimes used disagree in magnitude on the engineered
partition, so only the direction of that effect is treated as robust.
Externally, the study uses a single low-resolution dataset and mixes two
federation sizes, so no cross-scale or cross-dataset generalisation is claimed.
At the construct level, the rare-class set includes mid-support melanoma on
clinical grounds, which is a deliberate choice rather than a purely statistical
one.
% Rest: end with an explicit "what cannot be concluded" sentence (no universal
%   ranking; no clinical validation).
% Transition: "We summarise and answer the research questions in Chapter 7."


%% ----------------------------------------------------------------------------
\section*{7 Conclusion (opening sketch)}
This thesis set out to determine when, rather than whether, a federated optimiser
preserves rare-class performance on imbalanced medical data. Across statistical
heterogeneity, straggler asymmetry, learning-rate asymmetry and loss-side
correction, the results converge on one rule: rare-class performance is retained
precisely when the rare-client update continues to reach the global model. The
optimisers earn their place not as universal winners but as instruments that, in
specific regimes, protect or forfeit that signal.
% Rest: 7.2 answers RQ1-RQ4 in order; 7.3 restates contributions; 7.4
%   recommendations by regime; 7.5 future work.


%% ----------------------------------------------------------------------------
\subsection*{7.5 Future work (model paragraph)}
Several extensions follow naturally. Broadening the optimiser set to variance-
reduced and adaptive methods such as SCAFFOLD \citep{karimireddy2020scaffold} and
adaptive federated optimisers \citep{reddi2021adaptive} would test whether the
signal-flow account generalises beyond the three methods studied here.
Replicating the regime map on higher-resolution and additional medical benchmarks
would address the single-dataset limitation, and replacing simulated stragglers
with wall-clock deadlines would strengthen the system-heterogeneity claims.
Finally, the two interventions that act at different pipeline stages --
gamma-inexact aggregation and a weighted-CE local objective -- were tested only
separately; combining them is a concrete next step.
% Transition: closes the thesis.
